\documentclass{syllabus}

\begin{document}

\thispagestyle{fancy}
\coursecode{PHY 3960}
\renewcommand{\lms}{Brightspace}
\date{Fall 2026}
\author{}
\title{\coursecode: Portfolio Seminar B \\ {\large \mdseries Professional development as a physics major}}
\maketitle

\meeting{Tuesday: & 8:00 AM - 9:40 AM, DSC 054}
\officehours{MWF: & 2:00-3:30 PM}
\prerequisites{PHY 2970: & (Portfolio A)}


\section*{General Information}

\infotableport

\medskip
\noindent This seminar meets in the same period and room as PHY 4960 (Portfolio Seminar D); the two sections share their class sessions.

\section*{Course overview}

The second course in the Physics Portfolio Seminar sequence is focused on the personal and professional development of third year physics majors. In this course, students will be given opportunities to explore various possible career opportunities and reflect on their growth and accomplishments. Additionally, students will begin to build their portfolio by identifying experiences they hope to pursue before graduation and propose a self-directed project for the spring semester.

\subsection*{Student learning outcomes}
This course is structured to give students space to develop into better people primarily in the following two traits:
\begin{enumerate}[label=\textcolor{CarthageDark}{\bfseries\arabic*.}]
    \item\textbf{Magnanimity:} Aspiring to great deeds in a way that balances lofty goals and realism about personal capabilities.
    \item\textbf{Diligence:} Pursuing consistently the personal change and improvement necessary to achieve great deeds regardless of external incentives.
\end{enumerate}
Although the course will be organized so as to encourage this development, you are the only one who can determine whether to do the developing or not. In light of this, course assessment will be primarily structured around the following outcomes.

Upon completion of this course, a student who engages thoughtfully and earnestly with the course activities will be able to:
\begin{enumerate}[label=\textcolor{CarthageDark}{\bfseries\arabic*.}]
    \item Be aware of traditional and nontraditional physics career paths that contribute toward these goals.
    \item Explain how their personal values can give meaning to a variety of careers after graduation.
    \item Accurately assess the state of their knowledge and skills with respect to potential career paths and create a plan for how to develop knowledge and skill commensurate to career goals.
    \item Communicate effectively and professionally in both technical and non-technical terms while justifying physics-based decision making.
\end{enumerate}



\section*{Course components}

\subsection{Activities}
The lessons in the first portion of the course (joint with portfolio D students) will involve various activities. These might be reflections, discussions on readings, or updating your LinkedIn profile. The various week to week activities will make up one portion of your grade.

\textbf{AI usage is strictly prohibited on reflection assignments.} Reflections will be hand-written and require signing the following statement on each assignment:
\begin{quote}
    I acknowledge that the use of generative AI on this assignment represents a cowardly avoidance of engagement with my own personal experience.
\end{quote}

\subsection*{Participation} This course will involve frequent class discussions as well as student presentations; students will receive a participation score for each class session. A portion of this grade awarded based on attendance, with tardiness (arriving after 8:00) also being penalized. The remainder of the participation grade is awarded based on engagement with class activities (discussions or otherwise).

\subsection{Professional development plan} Over the course of the semester, there will be several assignments working toward a professional development plan. This will involve
\begin{enumerate}
    \item Identification of possible career trajectories and your motivations for pursuing these careers.
    \item Identification of skills or experience needed to succeed in the identified career trajectories.
    \item Meeting with a mentor to discuss a draft of your development plan.
    \item Proposal of a timeline of concrete activities on and off campus for developing these skills or obtaining these experiences over the remainder of their undergraduate career.
\end{enumerate}
Students are encouraged to have frequent discussions with their advisor or other mentors on this topic.



\subsection{Physics project proposal}
Portfolio C (in the spring semester) is an open-ended project-based course. Depending on how well-scoped and executed this project is, it could be a flop or a great addition to your portfolio. Based on skills you want to develop, you will be proposing (as a group) a project for next semester, which will allow us to hit the ground running on the project itself in the spring. Deliverables will include:
\begin{enumerate}
    \item Initial brainstorming of ideas, mapped to identified experience gaps.
    \item Decision matrix narrowing this list to three main contenders, considering a variety of stakeholder impacts.
    \item Preliminary investigation into all three finalist ideas, and rationale for selecting the final proposed idea.
    \item A final proposal for activities in the spring semester.
\end{enumerate}
Each stage will be accompanied by a presentation, with all except the narrowing step accompanied by reports.

\section*{Grading Scheme}
\subsection*{Component Weights}

\noindent\begin{tabular}{ll}
\textbf{Component} & \textbf{Weight}       \\ \hline
\textbf{Activities}            & 20\%       \\
\textbf{Participation}            & 10\%  \\
\textbf{Prof. dev. plan}            &     40\%   \\
\textbf{Project proposal}            & 30\%       \\
\end{tabular}%



\lettergrades{a) authentic engagement with class discussions and reflection assignments, and b) pattern of seeking help in developing a professional development plan from the instructor, academic advisor, or other identified mentor.}

\section*{Policies}
\subsection{General course policies}
Expect a response within one business day for emails sent to the instructor. To ensure a prompt response, format the subject line as: [\coursecode] *insert subject here*

Students are expected to arrive and leave on time. Cf. participation section for more information. I will attempt to arrive five minutes early and be available for at least five minutes after class.

\collegepolicies[College policies]

\end{document}
